\SourceLeaf{81}
\noindent oder auch

\[
\begin{aligned}fx^\lambda ={}& \left[Q^{-\frac{\lambda\nu}{n}}\right]_{\text{Coeff}\,q^{-\nu\lambda}}\\ &{}+ \frac{\lambda u}{\lambda+1}\left[Q^{-\frac{(\lambda+1)\nu}{n}}\right]_{\text{Coeff}\,q^{-\nu\lambda}}\\ &{}+ \frac{\lambda uu}{\lambda+2}\left[Q^{-\frac{(\lambda+2)\nu}{n}}\right]_{\text{Coeff}\,q^{-\nu\lambda}}+.\,.\,.\,..
\end{aligned}
\]

\noindent Setzt man demnach $\lambda = \dfrac{\theta}{\nu}$, so erhält man

\[
\begin{aligned}q^\theta ={}& u^{\frac{\theta}{\nu}}\left[Q^{-\frac{\theta}{n}}\right]_{\text{Coeff}\,q^{-\theta}}\\ &{}+ \frac{\theta}{\theta+\nu} .  u^{\frac{\theta+\nu}{\nu}}\left[Q^{-\frac{\theta+\nu}{n}}\right]_{\text{Coeff}\,q^{-\theta}}\\ &{}+ \frac{\theta}{\theta+2\nu} .  u^{\frac{\theta+2\nu}{\nu}}\left[Q^{-\frac{\theta+2\nu}{n}}\right]_{\text{Coeff}\,q^{-\theta}}+.\,.\,.\,.
\end{aligned}
\]

\noindent oder

\[
\begin{aligned}q^\theta ={}& p^{\frac{\theta m}{n}}\left[Q^{-\frac{\theta}{n}}\right]_{\text{Coeff}\,q^{-\theta}}\\ &{}+ \frac{\theta}{\theta+\nu} .  p^{\frac{(\theta+\nu)m}{n}}\left[Q^{-\frac{\theta+\nu}{n}}\right]_{\text{Coeff}\,q^{-\theta}}\\ &{}+ \frac{\theta}{\theta+2\nu} .  p^{\frac{(\theta+2\nu)m}{n}}\left[Q^{-\frac{\theta+2\nu}{n}}\right]_{\text{Coeff}\,q^{-\theta}}+.\,.\,.\,..
\end{aligned}
\]

\Needspace{7\baselineskip}Setzt man also $q^\theta = P$, so gilt

\[
[P]_{\text{Coeff}\,p^\mu} = \frac{\theta m}{\mu n}\left[Q^{-\frac{\mu}{m}}\right]_{\text{Coeff}\,q^{-\theta}}.
\]

\noindent Man kann dieses Ergebnis auch so ausdrücken: $\dfrac{dP}{dp}$ ist derjenige Term in der Entwicklung von

\[
\frac{m\theta\left(\dfrac{p^m}{Q}\right)^{\frac{\theta}{n}}q^\theta}{np\left(1-\left(\dfrac{p^m}{Q}\right)^{\frac{\nu}{n}}\right)}
\]

\noindent nach Potenzen von $q$, der kein $q$ enthält.

Daraus ergäbe sich folgende allgemeine Aufgabe. Die Beziehung zwischen $p$ und $q$ ist durch die Gleichung gegeben:

\[
P = Q,
\]

\noindent wobei $P$ eine Funktion von $p$ und $Q$ eine Funktion von $q$ von der Form

\[
q^n+aq^{n+\nu}+bq^{n+2\nu}+cq^{n+3\nu}+\text{etc.}
\]

\noindent ist. Gesucht ist die Darstellung von $fq$ durch eine Funktion von $p$.

